\chapter{KẾT LUẬN VÀ HƯỚNG PHÁT TRIỂN}
\label{chap:conclusion}

\section{Kết luận}

Nhóm đã hoàn thành quy trình khai phá dữ liệu hoàn chỉnh trên bộ dữ liệu giọng nói
Parkinson của UCI, với ba câu hỏi nghiên cứu được trả lời như sau.

\textbf{RQ1 -- Chất lượng dữ liệu.} Bộ dữ liệu canonical gồm 195 bản ghi $\times$ 23 cột,
không có giá trị thiếu hay dòng trùng lặp. Nhãn mất cân bằng (147 Parkinson / 48 khỏe),
thang đo đặc trưng chênh lệch lớn, 19/22 feature có candidate outlier theo IQR và 47/195
mẫu chứa ít nhất một candidate outlier; 14/22 feature có độ lệch $|\text{skew}| > 1$. Từ
đó nhóm áp dụng chuẩn hoá z-score và kiểm định chéo phân tầng để giảm ảnh hưởng của hai
đặc điểm này.

\textbf{RQ2 -- Phân loại có giám sát.} Trên 10-fold CV, ba mô hình đạt F-score macro 0,8757
(KNN), 0,8904 (SVM) và 0,8737 (Random Forest). SVM có kết quả cao nhất nhưng chênh lệch
giữa các mô hình nhỏ so với độ biến động giữa các fold; Random Forest cho thấy dấu hiệu
quá khớp (F-score huấn luyện 1,0000). Kết quả cho thấy tín hiệu giọng nói mang thông tin
phân biệt bệnh lý ở mức tốt.

\textbf{RQ3 -- Phân cụm không giám sát.} K-Means chọn $K=2$ theo silhouette (0,4373) nhưng
ARI chỉ $-0{,}0978$; DBSCAN đạt $\varepsilon = 3{,}37$, \texttt{min\_samples} = 5 cho một
cụm lớn, 12 điểm nhiễu và ARI $-0{,}0639$. Cấu trúc cụm trong không gian 22 đặc trưng do
đó không khớp ranh giới bệnh lý \texttt{status}; phân cụm đóng vai trò đối chiếu, không thay
thế mô hình có giám sát.

Ngoài kết quả kỹ thuật, đề tài còn rèn được quy trình làm việc nhóm: quản lý mã nguồn theo
nhánh, review từng pull request, đóng băng đầu vào chuẩn và chốt hai cổng GATE để đảm bảo
mọi kết quả nhất quán, tái lập được.

\section{Hạn chế}

\begin{itemize}
    \item Mẫu nhỏ (195 bản ghi, khoảng 32 chủ thể theo mô tả UCI) và mất cân bằng nhãn; mỗi
          fold kiểm tra chỉ khoảng 20 mẫu, trong đó khoảng 5 mẫu lớp khỏe nên điểm số giữa
          các fold dao động lớn.
    \item Cột \texttt{name} đã được loại khỏi canonical CSV nên không thực hiện kiểm định
          chéo theo nhóm chủ thể; bản ghi của cùng một người có thể rơi vào cả huấn luyện và
          kiểm tra, khiến điểm test có thể cao hơn thực tế.
    \item PCA hai chiều chỉ giữ 70,2\% phương sai; phân cụm không tách được cấu trúc khớp
          nhãn, DBSCAN chỉ tạo một cụm lớn với 12 điểm nhiễu.
    \item Mới so sánh ba mô hình với lưới tham số nhỏ; chưa tinh chỉnh ngưỡng quyết định,
          chưa huấn luyện trên tập dữ liệu độc lập.
    \item Một số thành viên chạy trên phiên bản thư viện khác nhau (ví dụ scikit-learn
          1.9.1); số liệu vẫn khớp khi kiểm tra chéo nhưng lần chạy cần ghi rõ môi trường.
\end{itemize}

\section{Hướng phát triển}

\begin{itemize}
    \item Bổ sung đặc trưng, lựa chọn đặc trưng (feature selection) và thử nghiệm thêm các
          mô hình hiện đại hơn; nghiên cứu kiểm định chéo theo nhóm chủ thể
          (\texttt{GroupKFold}) để khắc phục rò rỉ theo người.
    \item Điều chỉnh đơn vị/nhãn lưỡng phân hóa (threshold, class-weight) phù hợp với mục
          tiêu sàng lọc; đánh giá thêm độ tin cậy (calibration) và toàn diện (độ nhạy, độ
          đặc hiệu).
    \item Kết hợp nhãn với phân cụm (semi-supervised) hoặc học sâu trên tín hiệu âm thanh
          thô, và xây dựng giao diện nhỏ minh hoạ pipeline cho người dùng.
    \item Hoàn thiện hạ tầng tái lập: khóa phiên bản thư viện, thêm kiểm tra tự động cho
          dữ liệu đầu vào và đầu ra, chạy pipeline bằng một lệnh duy nhất.
\end{itemize}

\section{Đóng góp của nhóm và nghiệm thu}

Toàn bộ 21 nhiệm vụ trong bảng theo dõi đã được thực hiện và nghiệm thu \textbf{Done};
hai cổng GATE \#1 và GATE \#2 đã ký. Các pull request đã merge gồm \#5, \#7, \#8, \#13,
\#14, \#16 và \#17; mỗi nhiệm vụ chỉ được đóng sau khi có review độc lập.

\begin{table}[H]
\centering
\caption{Trạng thái nghiệm thu của từng thành viên}
\begin{tabular}{|c|l|l|}
\hline
\textbf{Mã} & \textbf{Thành viên} & \textbf{Trạng thái / bằng chứng} \\
\hline
M01 & Đặng Hoàng Ân   & Hoàn thành; ký GATE \#1, GATE \#2 \\
\hline
M02 & Lê Thiên Lộc    & Hoàn thành; PR \#5 đạt review \\
\hline
M03 & Huỳnh Gia Hợp    & Hoàn thành; PR \#13 đạt review \\
\hline
M04 & Trần Hoàng Thông & Hoàn thành; PR \#7 đã merge \\
\hline
M05 & Ngô Thị Anh Thư   & Hoàn thành; PR \#8 đạt review \\
\hline
M06 & Nguyễn Hoàng Bách & Hoàn thành; PR \#14, \#16 đạt review \\
\hline
\end{tabular}
\end{table}

Nhóm hy vọng báo cáo đáp ứng yêu cầu môn Khai phá dữ liệu; mọi thiếu sót còn lại nhóm sẽ
bổ sung theo hướng dẫn của giảng viên.
